% -----------------------------------------------------------
\section{Heart}
% -----------------------------------------------------------
	\label{HEART}
	
% % ----------------------
% \subsection{See also}
% % ----------------------

% ----------------------
\subsection{1828 American Dictionary of the English Language} % *1828
% ----------------------
	\label{HEART-1828}

	\begin{compactitem}
		\item[1] A muscular viscus, which is the primary organ of the blood's motion in an animal body, situated in the thorax. From this organ all the arteries arise, and in it all the veins terminate. By its alternate dilatation and contraction, the blood is received from the veins, and returned through the arteries, by which means the circulation is carried on and life preserved.
		\item[2] The inner part of any thing; the middle part or interior; as the \textit{heart} of a country, kingdom or empire; the \textit{heart} of a town; the \textit{heart} of a tree.
		\item[3] The chief part; the vital part; the vigorous or efficacious part.
		\item[4] The seat of the affections and passions, as of love, joy, grief, enmity, courage, pleasure etc. The \textit{heart} is deceitful above all things. Every imagination of the thoughts of the \textit{heart} is evil continually. We read of an honest and good \textit{heart} and an evil \textit{heart} of unbelief, a willing \textit{heart} a heavy \textit{heart} sorrow of \textit{heart} a hard \textit{heart} a proud \textit{heart} a pure \textit{heart} The \textit{heart} faints in adversity, or under discouragement, that is, courage fails; the \textit{heart} is deceived, enlarged, reproved, lifted up, fixed, established, moved, etc.
		\item[5] By a metonymy, \textit{heart} is used for an affection or passion, and particularly for love.
		\item[6] The seat of the understanding; as an understanding \textit{heart}.
		\item[7] The seat of the will; hence, secret purposes, intentions or designs. There are many devices in a man's \textit{heart}. The \textit{heart} of kings is unsearchable. The Lord tries and searches the \textit{heart} David had it in his \textit{heart} to build a house of rest for the ark.
		\item[8] Person; character; used with respect to courage or kindess. 
		\item[9] Courage; spirit; as, to take \textit{heart}; to give \textit{heart}; to recover \textit{heart}.
		\item[10] Secret thoughts; recesses of the mind.
		\item[11] Disposition of mind.
		\item[12] Secret meaning; real intention.
		\item[13] Conscience, or sense of good or ill.
		\item[14] Strength; power of producing; vigor; fertility. Keep the land in \textit{heart}.
		\item[15] The utmost degree.
	\end{compactitem}

% ----------------------
\subsection{Oxford English Dictionary} % *OED
% ----------------------
	\label{HEART-OED}
	
	There are more definitions for ``heart'' in the OED; this is just a selection.

	\begin{compactitem}
		\item[I.2] The heart considered as the centre of the vital or life-maintaining functions of the body; the seat of life; the vital part or principle; (in some contexts) life itself.
		\item[II.5] The bodily organ considered or imagined as the seat of feeling, understanding, and thought.
		\item[II.5.a] In the most general sense: the mind (including the functions of feeling and volition as well as intellect).
		\item[II.6.a] The seat or repository of a person's inmost thoughts, feelings, inclinations, etc.; a person's inmost being; the depths of the soul; the soul, the spirit.
		\item[II.7] (One's) intent, will, purpose; inclination, desire.
		\item[II.9.a] The heart considered as the seat of the emotions generally; a person's emotional nature (often contrasted with the intellectual or rationalizing nature located in the \textit{head}.
		\item[II.11.a] The heart considered as the seat of courage or morale.
		\item[II.12] The heart considered as the seat of perception or of the intellectual faculties generally; (also) understanding, intellect, mind; (in later use \textit{esp}.) memory.
	\end{compactitem}

% ----------------------
\subsection{Scriptural Usage} % *SCRIPTURES
% ----------------------
	\label{HEART-SCRIPTURES}

	% -----
	\subsubsection{Old Testament} % *OT
	% -----
		\label{HEART-OT}
	
		% --
		\paragraph{KJV}
		% --
			\label{HEART-OT-KJV}
	
			\begin{tabular}{ll}
				Total occurrences in the OT & 777
			\end{tabular}
			
			% \begin{compactdesc}
				% \item[]
			% \end{compactdesc}
			
		% --
		\paragraph{Hebrew Bible}
		% --
			\label{HEART-HEBREW}
			
			DCH and TDOT both say that there is no difference in meaning between \he{leb} and \he{lebAb} (for DCH, see below; see also TDOT 7:407, which says ``There is no discernible semantic difference; on the contrary, \he{leb} and \he{lebAb} appear to be totally synonymous and interchangeable'').
		
			%
			\subparagraph{\texorpdfstring{\he{leb}}{}} % לֵב
			%
				\label{LEB} % whatever is in step bible without periods
				
				The overwhelming majority of instances of ``heart'' in the KJV are translations of \he{leb}.
			
				\begin{compactdesc}
					\item[HALOT 513b, PDF 569] the quivering, pumping organ, the \textbf{heart} % Use the 1994 PDF
						\begin{compactitem}
							\item[1] the organ
							\item[2] seat of vital force (\hyperref[PSALMS-22-26]{Psalms 22:27} [in the KJV this is actually verse 26]), of illness (\hyperref[ISAIAH-1-5]{Isaiah 1:5})
							\item[3] \textbf{one's inner self}, seat of feeling and emotions
							\item[4] \textbf{inclination, disposition}
							\item[5] \textbf{determination, courage}
							\item[6] \textbf{will, intention}
							\item[7] \textbf{attention, consideration, reason}
							\item[8] mind in general and as a whole
							\item[9] \textbf{conscience}
							\item[10] metaphorical \textbf{inside, middle}
							\item[11] the organised \textbf{strength} of \he{nEpE+s}
							\item[12] \textbf{God's heart}
							\item[13] miscellaneous
						\end{compactitem}
					\item[BDB 524a, PDF 543] inner man, mind, will, heart
						\begin{compactitem}
							\item \textit{inner part, midst}
							\item[I] seldom of things
							\item[II] elsewhere of men
							\item[II.1] \textit{the inner man} in contrast with outer
							\item[II.2] \textit{the inner man}, indefinite, \textit{soul}, comprehending mind, affections and will, with occasional emphasis of one or the other by means of certain verbs
							\item[II.3] specific reference to \textit{mind}
							\item[II.3.a] of \textit{one's own mind}, \textit{destitute of mind}, etc.
							\item[II.3.b] \textit{knowledge}, \textit{wise mind}
							\item[II.3.c] \textit{thinking, reflection}
							\item[II.4] specifically referring to \textit{inclinations, resolutions and determinations of the will}
							\item[II.5] specifically referring to \textit{conscience}
							\item[II.6] specifically referring to \textit{moral character}
							\item[II.7] for \textit{the man himself}
							\item[II.8] as \textit{seat of appetites}
							\item[II.9] as \textit{seat of the emotions and passions}
							\item[II.10] \textit{seat of courage}
						\end{compactitem}
					\item[DCH PDF 2056] \textbf{heart}. Distinctions among [the following] meanings often unclear; a physical sense (\#4) is often possible in \#1--\#3. \he{leb} with possessive suffix (especially \he{lb*y}) may often refer to the whole person. Apparently not different in meaning from \he{lebAb}.
						\begin{compactitem}
							\item[1] \textbf{mind, thinking, intention, understanding}, in reference to primarily intellectual processes of human beings or of cities, nations, etc. spoken of in human terms
							\item[1.a] thought, reason, knowledge, counsel, of human beings
							\item[1.b] wisdom, common sense
							\item[1.c] attention, memory, of human beings
							\item[1.d] ability, skill
							\item[2] \textbf{feelings} of human beings. Especially as sensitive to:
							\item[2.a] joy
							\item[2.b] pain, grief, sadness
							\item[2.c] weakness, anxiety, feart
							\item[2.d] courage, strength, of human beings
							\item[2.e] irritation, anger
							\item[2.f] contempt
							\item[2.g] jealousy
							\item[2.h] conscience
							\item[3] \textbf{will, inclination, disposition, personality}, of human beings
							\item[4] physical \textbf{heart, chest} of human being
							\item[5] \textbf{middle, depth, height}
						\end{compactitem}
				\end{compactdesc}
				
				% \begin{compactdesc}
					% \item[Some OT uses] \textbf{}
						% \begin{compactdesc}
							% \item[]
						% \end{compactdesc}
				% \end{compactdesc}
			
		% % --
		% \paragraph{Septuagint}
		% % --
			% \label{HEART-LXX}
		
			% %
			% \subparagraph{\texorpdfstring{\gr{sample word}}{}}
			% %
		
	% -----
	\subsubsection{New Testament} % *NT
	% -----
		\label{HEART-NT}
	
		% --
		\paragraph{KJV}
		% --
			\label{HEART-NT-KJV}
			
	
			% \begin{tabular}{ll}
				% Total occurrences in the NT & 
			% \end{tabular}
			
			\begin{compactdesc}
				\item[{\hyperref[MATTHEW-15-18]{Matthew 15:18--19}}] (2)
			\end{compactdesc}
			
		% % --
		% \paragraph{Greek New Testament} % *GREEK
		% % --
			% \label{HEART-GREEK}
		
			% %
			% \subparagraph{\texorpdfstring{\gr{sample word}}{}}
			% %
				% \label{KATALLAGE}
			
				% \begin{compactdesc}
					% \item[BDAG , PDF ] \textbf{}
						% \begin{compactitem}
							% \item 
						% \end{compactitem}
					% \item[CGL , PDF ] \textbf{}
						% \begin{compactitem}
							% \item 
						% \end{compactitem}
					% \item[LSJ , PDF ] \textbf{}
						% \begin{compactitem}
							% \item 
						% \end{compactitem}
				% \end{compactdesc}
				
				% \begin{compactdesc}
					% \item[Some NT uses] \textbf{}
						% \begin{compactdesc}
							% \item[]
						% \end{compactdesc}
				% \end{compactdesc}
		
	% % -----
	% \subsubsection{Book of Mormon} % *BOM
	% % -----
		% \label{HEART-BOM}
	
		% \begin{tabular}{ll}
			% Total occurrences in the BOM & 
		% \end{tabular}
		
		% \begin{compactdesc}
			% \item[]
		% \end{compactdesc}
		
	% % -----
	% \subsubsection{Doctrine and Covenants} % *DC
	% % -----
		% \label{HEART-DC}
	
		% \begin{tabular}{ll}
			% Total occurrences in the DC & 
		% \end{tabular}
		
		% \begin{compactdesc}
			% \item[]
		% \end{compactdesc}
		
	% % -----
	% \subsubsection{Pearl of Great Price} % *PGP
	% % -----
		% \label{HEART-PGP}
	
		% \begin{tabular}{ll}
			% Total occurrences in the PGP & 
		% \end{tabular}
		
		% \begin{compactdesc}
			% \item[]
		% \end{compactdesc}
		
% % ----------------------
% \subsection{Key Phrases} % *KEYPHRASES *PHRASES
% % ----------------------
	% \label{HEART-PHRASES}

	% \begin{compactdesc}
		% \item[] \textbf{\#times} | list
			% % \begin{compactdesc}
				% % \item[Bible anchor verses] \phantom{.}
					% % \begin{compactdesc}
						% % \item[Romans 5:5] \gr{}
					% % \end{compactdesc}
				% % \item[Commentary] ---
			% % \end{compactdesc}
	% \end{compactdesc}
	
% % ----------------------
% \subsection{Bible Dictionaries} % *DICTIONARIES
% % ----------------------
	% \label{HEART-BD}

% % ----------------------
% \subsection{Restoration Usage} % *RESTORATION
% % ----------------------
	% \label{HEART-RESTORATION}

	% % -----
	% \subsubsection{Joseph Smith} % *JS
	% % -----
		% \label{HEART-JS}
	
		% \begin{compactdesc}
			% \item[{\hyperref[KEY]{Date/Source}}]
		% \end{compactdesc}
		
	% % -----
	% \subsubsection{Temple Ordinances}
	% % -----
		% \label{HEART-TEMPLE}
	
		% \begin{compactdesc}
			% \item[{\hyperref[KEY]{Date/Source}}]
		% \end{compactdesc}
	
	% % -----
	% \subsubsection{Brigham Young} % *BY
	% % -----
		% \label{HEART-BY}
	
		% \begin{compactdesc}
			% \item[{\hyperref[KEY]{Date/Source}}]
		% \end{compactdesc}
	
% ----------------------
\subsection{Commentary and Studies} % *COMMENTARY *STUDIES
% ----------------------
	\label{HEART-COMMENTARY}

	% % -----
	% \subsubsection{Key Points}
	% % -----
		% \label{HEART-KEYS}
	
		% \begin{compactitem}
			% \item
		% \end{compactitem}
		
	% -----
	\subsubsection{Reference information}
	% -----
		\label{HEART-reference}
	
		% -----
		\paragraph{\texorpdfstring{\he{leb}}{} in TLOT, PDF 823--824}
		% -----
			\label{HEART-TLOT}
			
			...Physical, psychological, and intellectual functions are attributed to the human \he{leb}...
			
			...The psychological aspect of the \he{leb} is evidenced in that it accommodates the most varied emotions...
			
			... The intellectual functions of the \he{leb} include, first of all, perception...Recognition or remembrance also occurs in the \he{leb}...Strictly intellectual capacities are also a matter of the \he{leb}: insight, the capacity to evaluate a matter critically, and juristic equilibrium. This aspect of the \he{leb} is especially significant for wisdom thought...The \he{leb} of the sage enables sound speech, he gains insight in the nature of the time and its events...
			
			...Finally, the \he{leb} is also the seat of the will and deliberation...
			
			...The \he{leb}, then, encompasses all dimensions of human existence...The expression \he{b:leb} (with a personal suffix) and a verb of thought or speech indicates thoughts that one keeps to oneself and does not wish to communicate. The dream too, which reveals the most hidden and inaccessible regions of a person, plays out in the \he{leb}...
			
			(Starting on PDF 825, TLOT discusses the theological significance of \he{leb}.)
			
		% -----
		\paragraph{\texorpdfstring{\he{leb}}{} in TDOT, 7:412, 419--420}
		% -----
			\label{HEART-TDOT}
			
			...The \he{leb} functions in all dimensions of human existence and is used as a term for all the aspects of a person: vital, affective, noetic, and voluntative. The prescientific anthropology of the OT does in fact view the human person as composite, althouh the analysis is not everywhere the same. This analysis is never the primary focus, since the OT texts seek to know men and women only as they stand before God...
			
			...The documents of the OT look on the human \he{leb} as the noetic seat of a wide range of activities. While the senses engage in ``sense perception''...within the \he{leb} take place intellectual visualization (cognition and memory), thought, understanding, and attention. Finally, wisdom is pictured as residing in the heart.
			
			Initially, cognition in the \he{leb} is related to sense perception: it is prior to seeing with the eyes and hearing with the ears, because it initiates the operation of the senses...Then it comes to mean the preservation and internalization of what the senses perceive for the purpose of making judgments and decisions. Thus cognition in the \he{leb} is always understood as a compact whole, in that it denotes the total noetic ability of an individual. Its opposite, therefore, is not a failure of sensory percetpion, but inattention, heedlessness, and confusion to the point of ethically negative duplicity. Perception and cognition in the \he{leb} are directed accordingly: initially toward concrete objects like the land of Canaan, parental instruction, the words of the wise; then increasingly (especially in the context of Deuteronomy) toward theological objects like Torah, God's promise, or signs and wonders. In prophet contexts, the object may be a vision, the word of God, or his grace and mercy...
			
			...If the \he{leb} lacks wisdom, it is in the grip of ``folly.'' Folly shows itself as heedlessness, willfulness, and failure to see the larger picture...Folly as the absence of \he{leb} leads ot disorientation and senselessness, adultery, and contempt for one's neighbor; in a youth, ``lack of heart'' manifests itself as stupidity, pleasure in worthless pursuits, and delight in foolishness. The \he{leb} of the fool loses control over the tongue and even goes so far as to deny God...
			
			...\hyperref[NUMBERS-15-39]{Numbers 15:39} commands the Israelites to let the tassels on their garments constantly remind them of the danger of being misled by the \he{leb}. A particularly extreme and dangerous form of deception by the \he{leb} is the phenomenon of false prophets, who prophesy out of their own \he{leb} and not as commanded by Yahweh...